\documentclass[10pt]{article}
\providecommand{\FIGDIR}{fig}
\newcommand{\DIGESTDATE}{21 September 2026}
% ===================== PM Research Watch - shared preamble =====================
\usepackage[T1]{fontenc}
\usepackage[utf8]{inputenc}
\usepackage{amsmath,amssymb}
\usepackage{tgpagella}
\usepackage{tgheros}
\renewcommand{\sfdefault}{qhv}
\renewcommand{\ttdefault}{lmtt}
\usepackage{microtype}
\usepackage[english]{babel}

\usepackage[a4paper,top=17mm,bottom=17mm,left=17mm,right=17mm,
            headsep=6mm,footskip=10mm]{geometry}
\usepackage{graphicx}
\usepackage{xcolor}
\usepackage{colortbl}
\usepackage{tikz}
\usetikzlibrary{calc}
\usepackage[most]{tcolorbox}
\usepackage{titlesec}
\usepackage{fancyhdr}
\usepackage{booktabs}
\usepackage{array}
\usepackage{longtable}
\usepackage{enumitem}
\usepackage{multicol}
\usepackage{ragged2e}
\usepackage[hidelinks]{hyperref}

% ------------------------------------------------------------------ palette
\definecolor{Ink}    {HTML}{14262E}
\definecolor{Deep}   {HTML}{0F4C5C}
\definecolor{Teal}   {HTML}{1F7A8C}
\definecolor{Sky}    {HTML}{4FA3C4}
\definecolor{Sage}   {HTML}{6FA287}
\definecolor{Amber}  {HTML}{E3A008}
\definecolor{Clay}   {HTML}{B5643C}
\definecolor{Coral}  {HTML}{D1495B}
\definecolor{Violet} {HTML}{6A5B9E}
\definecolor{Slate}  {HTML}{4A6572}
\definecolor{Paper}  {HTML}{FBF9F5}
\definecolor{Mist}   {HTML}{EEEAE2}
\definecolor{Rule}   {HTML}{D9D3C9}

\pagecolor{Paper}
\color{Ink}

\hypersetup{colorlinks=true, linkcolor=Deep, urlcolor=Teal, citecolor=Teal,
            pdftitle={PM Research Watch}, pdfauthor={Automated daily digest}}

% ------------------------------------------------------------------ headings
\titleformat{\section}
  {\sffamily\bfseries\large\color{Deep}}{}{0pt}
  {}[{\vspace{1.2mm}\color{Amber}\titlerule[1.1pt]}]
\titlespacing*{\section}{0pt}{5.5mm}{2.8mm}

\titleformat{\subsection}
  {\sffamily\bfseries\normalsize\color{Teal}}{}{0pt}{}
\titlespacing*{\subsection}{0pt}{3.6mm}{1.4mm}

\setlist[itemize]{leftmargin=4.2mm, itemsep=1.1mm, topsep=1.3mm, parsep=0pt,
                  label=\textcolor{Teal}{\footnotesize$\blacktriangleright$}}

% ------------------------------------------------------------------ header/footer
\pagestyle{fancy}
\fancyhf{}
\renewcommand{\headrulewidth}{0pt}
\renewcommand{\footrulewidth}{0pt}
\fancyhead[L]{\sffamily\scriptsize\color{Slate}PM RESEARCH WATCH\;\textcolor{Rule}{\textbar}\;\DIGESTDATE}
\fancyhead[R]{\sffamily\scriptsize\color{Slate}Particulate matter monitoring \& health effects}
\fancyfoot[C]{\sffamily\scriptsize\color{Slate}\thepage}

% ------------------------------------------------------------------ components
% metric tile
\newcommand{\metric}[3]{%
  \begin{tcolorbox}[enhanced, colback=white, colframe=Rule, boxrule=0.5pt,
    arc=1.6mm, left=2.2mm, right=2.2mm, top=1.8mm, bottom=1.8mm,
    borderline west={1.7mm}{0pt}{#3}, nobeforeafter, valign=center,
    height=17mm, width=\linewidth]
    \RaggedRight
    {\sffamily\bfseries\fontsize{17}{17}\selectfont\color{#3}#1}\\[0.5mm]
    {\sffamily\fontsize{6.5}{7.4}\selectfont\color{Slate}#2}
  \end{tcolorbox}}

% pull-quote / key insight box
\newtcolorbox{keybox}[1][]{enhanced, colback=Mist, colframe=Mist,
  arc=1.6mm, left=3.2mm, right=3.2mm, top=2.4mm, bottom=2.4mm,
  borderline west={1.7mm}{0pt}{Amber}, #1}

% subtopic band
\newcommand{\band}[2]{%
  \vspace{3.2mm}\noindent
  \begin{tcolorbox}[enhanced, colback=#2, colframe=#2, arc=1.2mm,
    left=2.6mm, right=2.6mm, top=1.1mm, bottom=1.1mm, boxrule=0pt,
    width=\linewidth, nobeforeafter]
    {\sffamily\bfseries\small\color{white}#1}
  \end{tcolorbox}\vspace{1.6mm}}

% one paper entry: {short cite}{title}{journal}{doi}{body}
\newcommand{\paperentry}[5]{%
  \par\vspace{1.9mm}\noindent
  {\sffamily\bfseries\footnotesize\color{Ink}#1}\;%
  {\sffamily\fontsize{7}{8}\selectfont\color{Slate}#3}\\[0.4mm]
  {\footnotesize\itshape\color{Deep}#2}\\[0.6mm]
  {\footnotesize #4}\\[0.3mm]
  {\fontsize{6.6}{7.6}\selectfont\color{Slate}\texttt{doi:}\,\href{https://doi.org/#5}{\texttt{\detokenize{#5}}}}
  \par}

% caption style
\newcommand{\figcap}[1]{\par\vspace{1.0mm}
  {\fontsize{7.2}{8.6}\selectfont\color{Slate}\textsf{\textbf{Fig.}}\; #1\par}\vspace{1.5mm}}

\providecommand{\permil}{\textperthousand}

\setlength{\parindent}{0pt}
\setlength{\parskip}{0pt}
\renewcommand{\arraystretch}{1.18}


\begin{document}

% ==================================================================== MASTHEAD
\thispagestyle{empty}
\begin{tikzpicture}[remember picture, overlay]
  \fill[Deep] (current page.north west) rectangle
      ($(current page.north east)+(0,-67mm)$);
  \fill[Amber] ($(current page.north west)+(0,-67mm)$) rectangle
      ($(current page.north east)+(0,-68.4mm)$);
  \foreach \x/\y/\r in {22/12/0.9, 41/26/1.5, 63/9/0.7, 88/31/1.1, 112/17/1.9,
                        138/29/0.8, 159/13/1.3, 178/54/1.0, 196/21/1.6,
                        30/44/0.6, 74/57/1.2, 125/48/0.9, 168/6/0.8, 191/59/0.7,
                        52/18/1.0, 100/60/0.6, 148/20/1.1, 66/49/1.4, 205/45/0.8}
    \fill[Sky, opacity=0.30] ($(current page.north west)+(\x mm,-\y mm)$) circle (\r mm);
\end{tikzpicture}

\vspace*{4mm}
{\sffamily\fontsize{7.6}{9}\selectfont\color{Sky}%
 AUTOMATED DAILY LITERATURE DIGEST\;\textbullet\;PhD RESEARCH WATCH\par}
\vspace{2.2mm}
{\sffamily\bfseries\fontsize{27}{29}\selectfont\color{white}%
 Particulate Matter\par}
{\sffamily\fontsize{27}{29}\selectfont\color{Sky}%
 Monitoring \& Health Effects\par}

\vspace{5.0mm}
{\sffamily\fontsize{8.4}{10}\selectfont\color{white}%
 Issue 2026-09-21\;\;\textcolor{Amber}{\textbar}\;\;Entry window 21 September 2026 (full day)\;\;%
 \textcolor{Amber}{\textbar}\;\;18 records screened in scope\par}
\vspace{18mm}

% -------------------------------------------------------------- metric strip
\noindent
\begin{minipage}[t]{0.190\linewidth}\metric{18}{RECORDS IN SCOPE\\(27 EXCLUDED)}{Deep}\end{minipage}\hfill
\begin{minipage}[t]{0.190\linewidth}\metric{8}{OF 10 SUBTOPIC\\BANDS POPULATED}{Teal}\end{minipage}\hfill
\begin{minipage}[t]{0.190\linewidth}\metric{2}{EFFECT ESTIMATES\\FROM ONE COHORT}{Sky}\end{minipage}\hfill
\begin{minipage}[t]{0.190\linewidth}\metric{3}{TIER-A\\RECORDS}{Amber}\end{minipage}\hfill
\begin{minipage}[t]{0.190\linewidth}\metric{5}{RECORDS CARRIED ON\\METADATA ALONE}{Clay}\end{minipage}

\vspace{6mm}

% -------------------------------------------------------------- signal of day
\section*{\sffamily\bfseries\color{Deep}Signal of the day}
\vspace{-1.5mm}{\color{Amber}\rule{\linewidth}{1.1pt}}\vspace{2.5mm}

\begin{keybox}
{\sffamily\bfseries\small\color{Deep}A consumer wearable, three sampling days per
pregnancy, and a 0.43\,SD deficit in femur length: the exposure instrument that made this
study possible is the same one whose calibration error is not reported.}\\[1.4mm]
\footnotesize
Tamire et al.\ (\textit{J Expo Sci Environ Epidemiol}) followed pregnant women in
\textbf{rural Butajira, Ethiopia} from the early second trimester, measuring personal
PM$_{2.5}$ over \textbf{24 hours with an Atmotube Pro} at roughly 20, 30 and 36 weeks and
tracking fetal growth by ultrasound at the same visits. An \textbf{IQR increase of
27.0\,$\mu$g\,m$^{-3}$} was associated with a \textbf{0.43\,SD decrease in femur length}
($\beta$ --0.43, 95\,\% CI --0.51 to --0.35) and a \textbf{0.17\,SD decrease in abdominal
circumference} (--0.17, --0.21 to --0.13), with \textbf{no significant association for head
circumference or biparietal diameter} under either concurrent or lagged exposure.
\textbf{The dissociation is the result, not the magnitude.} Skeletal and soft-tissue growth
respond while cranial growth is spared --- the brain-sparing pattern --- and that pattern is
hard to produce with a nutritional or socioeconomic confounder acting on growth as a whole.
Three limitations bite, and the first is the one that will be ignored. \textbf{The exposure
contrast is enormous by construction}: an IQR of 27\,$\mu$g\,m$^{-3}$ describes a
biomass-cooking population whose range does not overlap the cohorts these coefficients will
inevitably be compared against, so the $\beta$ is not transportable even if it is correct.
\textbf{Three 24-hour samples cannot represent a 16-week gestational window} --- but that
error is non-differential and biases toward the null, so it does not manufacture the femur
result. And the abstract does not report $n$.
\textbf{What makes this the signal of the day for a sensing programme is the instrument.}
This design --- repeated personal exposure across pregnancy in a rural low-income setting
--- was not affordable with gravimetric personal sampling, and it is now being run on
consumer-grade optical wearables. \textbf{The paper reports no sensor calibration, no
co-location against a reference, and no humidity correction.} Biomass smoke is the single
worst case for a nephelometric low-cost sensor: high relative humidity near a cookfire, a
fresh combustion aerosol whose size distribution and refractive index differ sharply from
the factory calibration aerosol. Any multiplicative bias in that instrument maps directly
onto the IQR and therefore onto $\beta$. The finding is plausible and the direction is
almost certainly right; the coefficient carries an uncertainty the paper does not quantify
and could have, with one co-location.\par
\end{keybox}

\vspace{3mm}

% -------------------------------------------------------------- what changed
\section*{\sffamily\bfseries\color{Deep}What changed today}
\vspace{-1.5mm}{\color{Amber}\rule{\linewidth}{1.1pt}}\vspace{2.5mm}

\footnotesize
\begin{itemize}

\item \textbf{Where you work is dirtier than
where you live, and the gap is socioeconomically patterned in the direction nobody assumes.}
A whole-population linkage of the \textbf{2022 Irish Census} to EPA pollutant surfaces
(\textit{Health Place}) assigned every person both a residential and a work/school address.
Mean annual PM$_{2.5}$, PM$_{10}$ and --- most sharply --- NO$_2$ are \textbf{higher at the
workplace than at home}, with O$_3$ running the other way, the expected NO$_x$-titration
signature. The contrast is \textbf{larger for more advantaged workers}, who commute from
cleaner suburbs into denser cores. Residence-only exposure assignment therefore does not
merely add noise; \textbf{it adds noise correlated with socioeconomic position}, which is
the one property an exposure proxy must not have in a study that adjusts for it. For
network design this is the strongest argument in this archive that \textbf{siting should
follow daytime population, not residential population}.

\item \textbf{HEPA filtration reliably cuts
indoor PM$_{2.5}$ and unreliably helps patients.} A synthesis of randomised trials in
asthma and COPD (\textit{Expert Rev Respir Med}) reports consistent, substantial exposure
reduction alongside \textbf{heterogeneous clinical effect} --- symptom, peak-flow and
quality-of-life gains concentrated in highly symptomatic or high-exposure subgroups, and
\textbf{most spirometric endpoints unmoved}. The authors' list of why is the contribution:
short follow-up, almost no adult studies, \textbf{device specifications unreported and
adherence unmeasured}. Those last two are measurement failures, not design failures, and a
logging sensor inside the intervention arm would close both.

\item \textbf{The instrumentation axis stayed
dark for a second day, but the deposits improved.} The \textbf{PubMed sensing query
returned zero again}, OpenAlex returned HTTP 429 and arXiv HTTP 406, so the instrumentation
legs rested entirely on Crossref-by-ISSN for the second consecutive issue. Unlike the 20th,
\textbf{the Copernicus (\textit{ACP}) and ACS (\textit{EST}) deposits carried full
abstracts}, which is why \textbf{metadata-only fell from 5 of 12 (42\,\%) to 5 of 18
(28\,\%)} without any change on the retrieval side. The five that remain are Elsevier and
Taylor \& Francis deposits, and they include the most directly relevant instrumentation
title of the window --- a method for measuring \textbf{penetration efficiency of low-cutoff
cyclones}, i.e.\ the calibration that defines what a size-selective inlet actually admits.
\textbf{The retry list now stands at 28 records.}

\item \textbf{The forest plot had a latent
ratio-scale assumption, and this issue is the first to trip it.} Every effect estimate this
watch has carried since July has been a ratio measure, so three assumptions were
load-bearing and invisible: a logarithmic axis, \textbf{a null line hardcoded at 1.0}, and a
geometric-mean span floor. The two Butajira estimates are \textbf{SD-scale regression
coefficients and both are negative}, and \texttt{(lo\,$\times$\,hi)$^{0.5}$} duly returned a
\textbf{complex number}, which killed the build in a comparison rather than at the axis.
Estimates are now \textbf{partitioned by scale}, each panel drawing its own axis, null and
tick ladder; a ratio metric whose lower bound is $\leq 0$ is routed to the difference panel
defensively. Worth recording that the colour map has had a \texttt{beta} entry since the
palette was written --- \textbf{the colour was anticipated, the axis was not.}

\item \textbf{Two upwind results argue the same
thing from opposite ends of India.} At \textbf{Sonipat}, the upwind receptor for Delhi-NCR
(\textit{ACP}), composition-based PM$_{2.5}$ exceeded \textbf{300\,$\mu$g\,m$^{-3}$} in two
post-monsoon haze episodes with organic aerosol at 65\,\% of non-refractory mass and
\textbf{57--60\,\% of it already oxygenated} --- the haze is aged secondary mass before it
reaches the city. In \textbf{Delhi} itself, replacing globally prescribed refractive indices
with local ones inside an operational forecast (\textit{ACP}, AIRWISE) cut surface
shortwave by up to \textbf{80\,W\,m$^{-2}$}, lowered boundary-layer height by
\textbf{200\,m} and \textbf{raised} PM$_{2.5}$. The feedback runs the wrong way for
forecasting: \textbf{getting the optics right makes the episode worse}, so a network
assimilating into a model with generic refractive indices inherits a low bias during exactly
the episodes that matter.

\item \textbf{A corpus defect was fixed ahead of
the monthly it would have corrupted.} \texttt{state/corpus/2026-08-03.json} carried a
\textbf{six-element} \texttt{LIFECOURSE} where the schema has five --- a trailing ``not
life-course specific'' overflow bucket holding 19 records. Because
\texttt{\_lifecourse\_range} sizes its accumulator from \textbf{the first issue it meets in
the range}, that record could silently widen every pooled life-course figure whose window
opened on it. The overflow bucket has been dropped and all 57 archived issues now conform.
One gap remains and is not a corruption: \texttt{2026-07-27} carries no
\texttt{LIFECOURSE} at all and is skipped cleanly by the rollup.

\end{itemize}

\vspace{2mm}

\begin{keybox}
{\sffamily\bfseries\footnotesize\color{Deep}Window continuity}\\[1.2mm]
{\fontsize{7.2}{8.8}\selectfont
The 20 September issue closed with \texttt{last\_entry\_date} at 20 September and explicitly
deferred the 21st, having run at 12:44 CDT with nine hours of that day's deposits
outstanding. This issue collects \textbf{21 September in full}, harvested at 23:15 CDT, so
the window is \textbf{continuous with the previous issue and leaves no gap}. The run crossed
midnight during the build. \textbf{22 September is not opened here}: it was minutes old when
this issue compiled, and closing it would both produce a near-empty issue and consume the
day. The next scheduled run harvests 22 September in full. The weekly rollup is not due ---
W38 shipped with the 19 September (Saturday) issue, per the standing Saturday cadence.\par}
\end{keybox}

\vspace{3mm}
% ==================================================================== ANALYTICS
\section*{\sffamily\bfseries\color{Deep}Corpus analytics}
\vspace{-1.5mm}{\color{Amber}\rule{\linewidth}{1.1pt}}\vspace{3mm}

\noindent\includegraphics[width=\linewidth]{\FIGDIR/f1_subtopics.png}
\figcap{Subtopic assignment of the 18 in-scope records; eight of the ten bands are
populated. The shape is the opposite of the 20 September issue's flat distribution:
\textbf{\textit{Exposure assessment \& modelling} (5) and \textit{Sensing, forecasting \&
instrumentation} (4) together take half the corpus}, which is the watch's own axis and the
first time this week it has led. Read that against the \textbf{four health-endpoint records
in the whole issue} --- one each in reproductive, respiratory, neuro and mechanistic --- and
the issue is better described as an instrumentation and exposure day with four clinical
records attached than as a balanced one. \textit{Cardiovascular \& metabolic} and
\textit{Other clinical endpoints} are empty; on 20 September the empty pair was
cardiovascular and neuro, so cardiovascular has now been absent two issues running.}

\vspace{2mm}
\noindent\includegraphics[width=\linewidth]{\FIGDIR/f2_design_endpoint.png}
\figcap{Left --- study architecture. \textbf{\textit{Metadata only} is again the largest
single slice at five of eighteen}, but the share has fallen from 42\,\% to \textbf{28\,\%}
and it no longer outnumbers characterised field work: \textbf{measurement campaigns take
four} and \textbf{chamber/laboratory three}, with single cohort, cross-sectional, acute,
modelling, synthesis and toxicology records behind them. Nine architectures across eighteen
records is a genuinely diverse day. Right --- health endpoint, canonicalised.
\textbf{``No health endpoint'' takes fourteen of eighteen}, the highest share this archive
has recorded, and unlike the 20th that is \textbf{mostly real rather than an artefact of
missing abstracts}: \textbf{nine of those fourteen} are exposure, instrument, forecast,
filtration or mitigation studies that genuinely measure no health outcome, against
\textbf{five} that are endpoint-blank only because their abstract is missing. On
20 September the same split was five real to five blank, so the growth in this bar is
real studies, not missing ones. The four records that do carry an endpoint are singletons
in four different categories: respiratory, reproductive, neuro/mental health and
\textit{in vitro} toxicology.}

\vspace{2mm}
\noindent\includegraphics[width=\linewidth]{\FIGDIR/f3_metric_geography.png}
\figcap{Left --- particle metric. The spread is wider than usual because four records
measure something other than PM mass: \textbf{equivalent black carbon}, \textbf{oxygenated
organic molecule fluxes}, \textbf{aerosol refractive index} and \textbf{respirable coal dust
$D_{50}$}. Two records measure PM$_{2.5}$ with a \textbf{consumer-grade optical wearable}
(Butajira) or a \textbf{minute-level mobile sensor} (Hong Kong) --- the first issue this
month in which two independent health studies both rest on low-cost sensing hardware. Right
--- geography. \textbf{China takes four} (national vertical profile, a central-China flux
site, a deep coal face, and Hong Kong) and \textbf{South Asia three}, all Indian: Sonipat,
Delhi and semi-arid western India. \textbf{The six in \textit{Global / multi-region} are
not a finding} --- five are non-geographic by construction (two bench filtration rigs, one
\textit{in vitro} study, two metadata-only records with no recoverable site) and the sixth
is a multi-country trial synthesis. Sub-Saharan Africa's single record is the tier-A one.}

\vspace{2mm}
\noindent\includegraphics[width=\linewidth]{\FIGDIR/f6_lifecourse.png}
\figcap{Life-course windows addressed; counts are non-exclusive and a record is counted
only where the study design is specific to a stage. \textbf{Two of the five windows are
populated} --- \textbf{preconception/in utero} by the Butajira pregnancy cohort and
\textbf{working-age adulthood} by three occupational or ambulatory studies (the Japanese
teaching laboratory, the Chinese coal face, the Hong Kong momentary-stress panel).
\textbf{Childhood, adolescence and older adults are all empty}, and the 65+ gap is the one
worth noting: it was populated on both the 19th and the 20th and is the stage where the
mortality literature concentrates. Thirteen records are not life-course specific and are
counted in no stage, which is why the strip sums to five rather than eighteen.}

\vspace{4mm}

\noindent\includegraphics[width=\linewidth]{\FIGDIR/f4_heatmap.png}
\figcap{Subtopic $\times$ study-architecture cross-tabulation. The matrix is sparse ---
\textbf{eighteen records spread over fourteen of the seventy-two available cells}, and
\textbf{none holds more than two}. \textbf{Four cells hold two}:
\textit{Exposure} $\times$ \textit{Metadata only}, \textit{Sensing} $\times$
\textit{Metadata only}, \textit{Exposure} $\times$ \textit{Measurement campaign} and
\textit{Occupational} $\times$ \textit{Chamber/laboratory}. The first two are the day's
structural story: \textbf{the unreadable records concentrate in precisely the two bands
this watch exists to track}, because those are the bands served by the Elsevier and
Taylor \& Francis deposits that carry no abstract. The characterised instrumentation work
in this issue is all Copernicus. Read the other two as the healthy pair --- exposure work
is being done by measurement campaigns and occupational work on the bench, which is what
those bands should look like.}

\vspace{2mm}
\noindent\includegraphics[width=\linewidth]{\FIGDIR/f5_forest.png}
\figcap{\textbf{Both extractable effect estimates, and the first difference-scale estimates
this watch has carried} --- hence a linear axis with the null at \textbf{zero}, not the
logarithmic ratio axis every previous issue has used. Both come from \textbf{one cohort}
(Butajira) and this is therefore \textbf{not a synthesis and must not be read as one}; it
is one study's two-coefficient growth matrix. The internal comparison is the informative
part: \textbf{femur length (--0.43\,SD) falls roughly two and a half times as far as
abdominal circumference (--0.17\,SD)} per the same IQR contrast, and the two intervals do
not overlap, so the skeletal-versus-soft-tissue ordering is internally significant rather
than an eyeball impression. Head circumference and biparietal diameter are \textbf{absent
from this figure because they were null}, not because they were unreported --- the
brain-sparing pattern is carried by what is missing here. The contrast is per study-specific
IQR (27.0\,$\mu$g\,m$^{-3}$), not per 10\,$\mu$g\,m$^{-3}$, so nothing here is comparable to
estimates in other issues.}

\vspace{4mm}

% ==================================================================== DIGEST
\section*{\sffamily\bfseries\color{Deep}Paper-level digest}
\vspace{-1.5mm}{\color{Amber}\rule{\linewidth}{1.1pt}}\vspace{1mm}

{\fontsize{7.4}{9}\selectfont\color{Slate}Ordered by cluster. Each entry gives the
methodological core and the finding that matters, not an abstract paraphrase. Records
marked \textit{metadata only} carry no recoverable abstract and assert no finding.\par}

\band{Neuro / mental health\hfill 1 record}{Deep}

\paperentry{Hong Kong momentary stress 2026}
{An Explainable Multi-Modal Deep Learning Framework for Investigating the Influences of Dynamic Environmental Exposure on Momentary Stress.}
{Environ Sci Technol}
{4,342 momentary stress reports against 125,440 minute-level environmental exposure records; the model reaches 78.26\% accuracy and 0.828 AUC-ROC. Two results are worth extracting from the architecture. First, exposure in the preceding 25 minutes dominates prediction, which is a direct empirical argument for minute-resolution mobile monitoring over daily or residential-annual assignment for acute endpoints. Second, the explainability analysis recovers thresholds rather than slopes - 25 $\mu$g\,m$^{-3}$ for PM$_{2.5}$, 55 dBA noise, 32 C, 40-70\% humidity, 25\% greenery visibility - consistent with nonlinearity that a linear exposure-response term would average away. The weaknesses are structural and serious: stress is self-reported into the same phone that carries the sensors, so reporting and exposure share a device and a context; the design is correlational with no within-person randomisation of exposure; the thresholds are read off a gradient-based attribution of a black-box model, not estimated with uncertainty, so they should be treated as hypotheses; and PM$_{2.5}$ co-varies with noise and temperature in ways a deep model can blend invisibly. Still the clearest current demonstration of what wearable-grade PM sensing buys over address-based assignment.}
{10.1021/acs.est.6c09466}

\band{Respiratory \& allergic\hfill 1 record}{Sky}

\paperentry{HEPA RCT synthesis 2026}
{HEPA air filtration for asthma and COPD: lessons from randomized controlled trials.}
{Expert Rev Respir Med}
{The consistent finding across trials is the one that should worry anyone who equates exposure reduction with clinical benefit: portable HEPA reliably and substantially lowers indoor PM$_{2.5}$, but clinical effects are heterogeneous. Several asthma trials improved symptoms, peak flow or quality of life - concentrated in highly symptomatic or higher-exposure subgroups - while most spirometric endpoints moved little or not at all. COPD evidence is thin and only suggestive. The authors attribute the gap to trial design rather than to absent efficacy, and list the specific failures: short follow-up, almost no adult studies, device specifications unreported, adherence unmeasured, no economic evaluation, and near-total absence of low-resource settings. That list is itself the contribution. Being a narrative rather than systematic synthesis with no pooled estimate and no formal risk-of-bias assessment, it cannot bound the effect, and the subgroup pattern it highlights is exactly the pattern selective reporting produces. For a sensing programme the actionable point is that unmeasured adherence and unreported device spec are measurement problems a logging sensor inside the intervention arm would largely solve.}
{10.1080/17476348.2026.2735502}

\band{Reproductive \& developmental\hfill 1 record}{Sage}

\paperentry{Tamire et al. 2026}
{Prenatal fine particulate matter (PM$_{2.5}$) exposure differentially impairs fetal skeletal and soft-tissue growth: a longitudinal study in rural Butajira, Ethiopia.}
{J Expo Sci Environ Epidemiol}
{Pregnant women recruited at or before 24 weeks, followed over three visits, March-June 2025. An IQR rise of 27.0 $\mu$g\,m$^{-3}$ in personal PM$_{2.5}$ was associated with a 0.43 SD fall in femur length (beta -0.43, 95\% CI -0.51 to -0.35) and a 0.17 SD fall in abdominal circumference (-0.17, -0.21 to -0.13); head circumference and biparietal diameter showed no significant association under either concurrent or lagged exposure. The dissociation is the finding - skeletal and soft-tissue growth respond while cranial growth is spared, which is the brain-sparing pattern and argues against a purely nutritional confounder. Three weaknesses bite. The exposure contrast is enormous by design (an IQR of 27 $\mu$g\,m$^{-3}$ implies a biomass-cooking population far outside the range of the cohorts these coefficients will be compared against), the sample is small and the paper does not report n in the abstract, and three 24-h samples cannot represent a 16-week gestational window - classical measurement error that, being non-differential, biases toward null and so does not explain the femur result. Directly relevant to distributed sensing: the exposure instrument is a consumer-grade optical wearable, and effect sizes this large resting on three wearable-days per pregnancy is exactly the design low-cost sensing makes newly affordable in LMIC settings - and exactly where sensor humidity and biomass-smoke calibration error, unreported here, would propagate straight into beta.}
{10.1038/s41370-026-00966-4}

\band{Mechanistic toxicology\hfill 1 record}{Violet}

\paperentry{CSE HASMC isoforms 2026}
{A comprehensive full-length transcriptome landscape of cigarette smoke-exposed HASMCs reveals extensive remodeling of mRNA isoforms and regulatory networks.}
{PLoS One}
{Oxford Nanopore long-read sequencing of human aortic smooth muscle cells at 0.6\% cigarette smoke extract, the midpoint of a biphasic viability curve (0.2\% mildly proliferative, 0.4-0.8\% progressively cytotoxic). 11,222 novel transcripts; 1,866 transcripts up in ferroptosis and cell division pathways, 493 down in oxidative phosphorylation and glutathione metabolism, 1,753 alternative polyadenylation events remodelling 3'UTRs, and 269 alternative splicing events hitting cell-cycle regulators and transcription factors including CREM, FOSL1 and HBP1. The claim is that high-dose cytotoxicity is executed post-transcriptionally rather than through differential gene expression alone, and the isoform resolution is what makes that visible - a short-read study would have reported the pathway shifts and missed the 3'UTR remodelling entirely. Admitted at tier C with an explicit caveat: cigarette smoke extract is not ambient PM, the cell line is not a vessel, a single 0.6\% dose is not a dose-response for the mechanism claimed, and no particle mass or size metric is reported, so nothing here is transferable to an ambient exposure-response coefficient. Carried for the ferroptosis and glutathione-depletion axis, which recurs in ambient PM vascular toxicology and is worth tracking.}
{10.1371/journal.pone.0354604}

\band{Exposure assessment \& modelling\hfill 5 records}{Clay}

\paperentry{Ireland Census linkage 2026}
{Ambient air pollution at home, work and school: Evidence from Ireland.}
{Health Place}
{Whole-population census microdata linked to Irish EPA pollutant surfaces at two addresses per person. Mean annual PM$_{2.5}$, PM$_{10}$ and - most sharply - NO$_2$ are higher at the work, school or college address than at the residence, while O$_3$ runs the other way, the expected NOx-titration signature of the same urban gradient. The socioeconomic result inverts the usual assumption: among workers the residential-to-workplace contrast is *larger* for more advantaged groups, because advantaged workers commute into dense urban cores from cleaner suburbs. Residence-only assignment therefore does not simply add noise, it adds noise correlated with socioeconomic position, which is the one thing an exposure proxy must not do in a study adjusting for it. Limits: both endpoints are modelled ambient surfaces at an address, so nothing here captures infiltration, indoor sources or time-activity weighting, and the paper measures the size of a bias without estimating its effect on any health coefficient. For network design this is the strongest argument yet that siting density should follow daytime population rather than residential population.}
{10.1016/j.healthplace.2026.103732}

\paperentry{Sonipat ACSM campaign 2026}
{Chemical characterization and source apportionment of carbonaceous aerosols during post-monsoon biomass burning and Diwali at an upwind site of Delhi.}
{Atmos Chem Phys}
{25 October to 15 November 2023, spanning two severe haze episodes and Diwali. Composition-based PM$_{2.5}$ exceeded 300 $\mu$g\,m$^{-3}$ in both episodes; organic aerosol was 65\% of non-refractory PM$_{2.5}$ with a daily mean near 140 $\mu$g\,m$^{-3}$, and oxygenated OA was 57-60\% of total OA, reaching 42.8 $\mu$g\,m$^{-3}$ of more-oxidised OOA at peak haze - so the haze is substantially aged secondary mass, not fresh plume. Biomass- or solid-fuel eBC was \textasciitilde{}78\% of campaign-integrated eBC. The authors are honest that this split is sensitive to the assumed absorption Angstrom exponent pair, which is the paper's real limitation: the headline 78\% is a function of a prior, and the sensitivity range is not propagated into the PMF factors. Single site, single season, three weeks. The result that matters upstream of a sensor network is that the mass is already loaded before the air reaches Delhi, so an urban-only network measures an urban increment on a large regional background and will misattribute regional transport to local sources unless it has an upwind anchor.}
{10.5194/acp-26-13227-2026}

\paperentry{W. India BC-dust 2026}
{Seasonal Variability and Interactions of Black Carbon and Mineral Dust over Semi-Arid Western India: Satellite and Ground-Based Observations.}
{Environ Sci Atmos}
{Only the opening of the abstract was deposited, so the design is legible but the numbers are not: the study characterises seasonal co-variation of anthropogenic black carbon and natural mineral dust in a semi-arid region where both are large, using satellite retrievals alongside ground-based observations. No effect sizes, retrieval-validation statistics or station counts are recoverable from the deposit, and nothing is claimed here beyond scope. Flagged for retry rather than summarised further. The reason to carry it at all is the mixing problem it addresses: BC and dust are the two components that most degrade low-cost optical PM sensing in arid regions, BC because it absorbs rather than scatters and dust because its coarse mode falls outside a nephelometric sensor's effective cut, and a seasonal climatology of their interaction is the calibration prior a regional network would need.}
{10.1039/d6ea00109b}

\paperentry{China aerosol vertical 2026}
{Spatiotemporal Dynamics and Vertical Distribution of Aerosols in China: Insights from Multi-Source Evidence from 2000 to 2022.}
{Atmos Pollut Res}
{\textsc{metadata only.} Metadata-only record: Elsevier deposited title, journal and DOI without an abstract, and no abstract was recoverable from Crossref, Europe PMC or PubMed. A 23-year national synthesis of aerosol vertical structure is potentially a tier-A input to satellite-to-surface PM$_{2.5}$ reconstruction - vertical profile shape is the dominant structural uncertainty when converting AOD to surface concentration - but nothing about method, platform mix or findings can be stated from the deposit, and inferring a design from a title is how this watch would start manufacturing results. Carried as a placeholder and added to the retry list; if the abstract appears it should be re-screened on its merits rather than admitted on the strength of its title.}
{10.1016/j.apr.2026.103198}

\paperentry{Athens NPF 2026}
{The influence of air mass origin and ammonia emissions on new particle formation at a Mediterranean urban site (Athens, Greece).}
{Atmos Res}
{\textsc{metadata only.} Metadata-only record; Elsevier deposit carried no abstract and none was recoverable. The topic is squarely in scope - ammonia-mediated new particle formation is the ultrafine source that governs urban particle *number*, and number is precisely the metric that mass-calibrated low-cost sensors do not measure and that regulatory PM$_{2.5}$ mass does not capture. Whether this paper quantifies nucleation rates, growth rates or only classifies event days cannot be determined from the deposit. Retry list.}
{10.1016/j.atmosres.2026.109334}

\band{Sensing, forecasting \& instrumentation\hfill 4 records}{Amber}

\paperentry{AIRWISE Delhi 2026}
{Improving aerosol-radiation interactions in the operational forecasting system - AIRWISE.}
{Atmos Chem Phys}
{Replacing globally prescribed refractive indices with Delhi-specific values inside the operational AIRWISE system cut surface downwelling shortwave by up to \textasciitilde{}80 W/m2 instantaneously (\textasciitilde{}43 W/m2 diurnally) during a severe post-monsoon episode, and monthly mean SWDOWN by 25-37 W/m2, a quarter to a third of total clear-sky aerosol-radiation interaction. The meteorological response is coherent and the right sign: surface temperature -0.2 C, near-surface wind -0.4 m/s, boundary layer height -200 m, daytime humidity +3-4\%, and consequently higher PM$_{2.5}$ - the shallower, calmer boundary layer traps the aerosol that caused the dimming. Validated against WiFEX surface radiation at Delhi airport with 62\% (December) and 35\% (January) bias reduction. Limits: one city, one season, refractive indices tuned where they are evaluated, and radiation bias reduction does not demonstrate PM$_{2.5}$ forecast skill improvement, which is the operational claim a reader will assume. Relevant because the feedback runs the wrong way for forecasting: getting optics right makes the pollution episode worse, so a network assimilating into a model with generic refractive indices inherits a systematic low bias during exactly the episodes that matter.}
{10.5194/acp-26-13189-2026}

\paperentry{Urban OOM flux 2026}
{Urban surface-atmosphere fluxes of selected pptv-level oxygenated organic molecules (OOMs) from eddy covariance observations.}
{Atmos Chem Phys}
{First urban eddy covariance measurement of OOM surface-atmosphere exchange at pptv level, and the contribution is as much methodological as observational - the paper works through how mass spectral integration time and water vapour variation corrupt flux retrieval at low signal-to-noise, which is the reason these fluxes had not been measured. Sixteen OOMs resolved into deposition-dominated, emission-dominated and bidirectional classes. Campaign-mean deposition was 4.6 umol/m2/d, 16.3\% of HNO$_3$ deposition and 2.2x the OOM emission flux, so the urban surface is a net sink. Formic acid dominated both directions (66.2\% of deposition, 84.9\% of emission), which means the 16-species budget is really a formic acid budget plus minor terms. One site, one month, hottest season - the seasonal extreme most favourable to the photochemistry being measured, so these are upper bounds. Bears on sensing indirectly but concretely: SOA formed from these precursors is secondary mass a low-cost optical sensor sees only after it has condensed, with a lag and a hygroscopicity that no dry-mass calibration captures.}
{10.5194/acp-26-13213-2026}

\paperentry{Low-cutoff cyclone method 2026}
{An improved method for measuring the penetration efficiency of low cutoff cyclones.}
{Aerosol Sci Technol}
{\textsc{metadata only.} Metadata-only record: Taylor \& Francis deposited no abstract, the landing page is bot-blocked, and neither Crossref nor a Consensus sweep of 2026 cyclone-characterisation work surfaced it. This is the most directly relevant instrumentation title of the window and the least legible. Low-cutoff cyclone penetration efficiency is the calibration that defines what a size-selective inlet actually admits, and an error there propagates into every downstream gravimetric or optical measurement behind that inlet - including the reference samplers that low-cost sensor networks are co-located against, where an inlet cut error becomes a systematic multiplicative bias in the correction factor rather than added noise. Top of the retry list; worth a manual full-text pull rather than waiting on the deposit.}
{10.1080/02786826.2026.2731025}

\paperentry{Medellin MABI eBC 2026}
{Multi-wavelength spectral attribution of equivalent black carbon in Medellin, Colombia (2010-2023): Disentangling fossil fuel and biomass burning using the MABI instrument.}
{Atmos Pollut Res}
{\textsc{metadata only.} Metadata-only record; no abstract deposited or recoverable. A 14-year multi-wavelength absorption record with fossil-versus-biomass attribution would be an unusually long instrument series for a Latin American city, and the same absorption-Angstrom-exponent dependence the Sonipat paper in this issue flags as its principal uncertainty applies here, so the two would be worth reading against each other once the text is available. Nothing further can be said without it; retry list.}
{10.1016/j.apr.2026.103203}

\band{Occupational \& indoor\hfill 3 records}{Amber}

\paperentry{JP laboratory exposure 2026}
{Assessment of personal exposure and spatial concentration variability in a small-scale chemical laboratory.}
{Ann Work Expo Health}
{Personal exposure concentrations ran consistently higher than co-located area measurements, mean values roughly an order of magnitude greater, reproduced across two independent measurement years. CFD then explains the gap mechanically: concentration varies strongly with height, so a fixed sampler placed away from the breathing zone of a person leaning over their own source cannot recover personal exposure, whatever its precision. The timing is regulatory - Japan makes personal exposure measurement mandatory from October 2026 - and the paper is effectively a demonstration that the outgoing area-based system was biased low by roughly 10x in this setting. Narrow and honest about it: one small teaching laboratory, student experiments rather than production work, passive samplers for vapour-phase analytes rather than PM. The 10x should not be transported to other geometries. It generalises to distributed PM sensing as a siting argument with a number attached - a fixed node samples the room, a person samples their own plume, and near a source the two differ by an order of magnitude, not a percentage.}
{10.1093/annweh/wxag082}

\paperentry{SAA wetting agent 2026}
{High-efficiency 'SAA' wetting agent for coal mine dust control: from preparation and mechanism to field application.}
{Environ Geochem Health}
{A three-component surfactant blend (sodium gluconate, AEO-3, AES) lowers surface tension and contact angle on coal, raises oxygen-containing surface functional groups, and agglomerates micron-scale dust - D50 rose 22.58\% over untreated. In a purpose-built simulation rig the suppression rate was 71.50\%, 2.17x conventional water mist, and the formulation was then taken underground on a real working face. The mechanism-to-field arc is the strength; the evaluation is the weakness. A suppression rate measured in the authors' own custom rig has no external calibration, the field test reports no control shift or personal sampling, and no respirable-fraction or silica-specific result is given, which is what determines pneumoconiosis risk - agglomerating dust into coarser particles raises D50 partly by removing mass from exactly the respirable size range the health endpoint cares about, so the health-relevant efficacy is almost certainly different from 71.5\% and is not reported. Relevant to occupational sensing as a reminder that a mass-based sensor reads an agglomeration intervention as a larger improvement than the lung sees.}
{10.1007/s10653-026-03491-3}

\paperentry{ZCND aerogel filter 2026}
{Biomimetic dual-network nanofibrous aerogels with robust interfacial assembly for multifunctional air filtration.}
{J Colloid Interface Sci}
{A zein/curcumin nanofibre plus down-fibre aerogel reached 99.51\%, 99.74\% and 99.93\% efficiency for PM$_{1}$, PM$_{2.5}$ and PM$_{10}$ at 27.39 Pa pressure drop, with 335.85 g/m2 dust-holding capacity, and - the result that distinguishes it - retained 98.67\% PM$_{1}$ efficiency after 12 h at 90\% relative humidity. It also adsorbs ammonia (104.95 cm3/g), shows visual NH$_3$ sensing, is antibacterial, and biodegrades in soil burial faster than commercial media. The target is livestock housing, where PM, NH$_3$ and bioaerosol co-occur and humidity is chronically high, so the multifunctionality is matched to a real setting rather than assembled for novelty. Standard caveats for this literature apply and are not minor: bench aerosol challenge under controlled conditions, no barn deployment, 12 h is not a service life, and nothing on mechanical durability under pulse cleaning or cost at scale. Noted here mainly because humidity-resistant PM$_{1}$ capture is the same physical problem, inverted, that defeats nephelometric low-cost sensors in barns - both fail when water condenses on the particle, one by fouling and one by over-reading.}
{10.1016/j.jcis.2026.141640}

\band{Burden, policy \& mitigation\hfill 2 records}{Coral}

\paperentry{N fertilizer NH$_3$-PM$_{2.5}$ 2026}
{Integrating nitrogen fertilizer management for mitigating agricultural ammonia emissions and PM$_{2.5}$ pollution.}
{Atmos Res}
{\textsc{metadata only.} Metadata-only record; no abstract deposited or recoverable from any leg. Agricultural ammonia is the control lever with the best cost-per-microgram in most secondary-inorganic-dominated airsheds and the one least often regulated, so a fertilizer-management framing is squarely in scope for the burden and mitigation track. Whether the paper is a modelling assessment, a field trial or a policy review is not determinable from the deposit. Retry list. Note it sits alongside the Athens new-particle-formation record in this same issue - two independent ammonia-centred papers deposited the same day, both illegible, which is a coverage gap rather than a coincidence worth reading into.}
{10.1016/j.atmosres.2026.109343}

\paperentry{PPS antifouling filter 2026}
{Robust and Self-Cleaning Fluorine-Free Superhydrophobic PPS Filters for High-Humidity Flue Gas Filtration.}
{Environ Sci Technol}
{Light-induced covalent grafting of a triple-cross-linked low-surface-energy network onto polyphenylene sulfide fibre membranes produces a fluorine-free antifouling bag-filter medium. Against the unmodified control it extends mean pulse-jet cleaning cycle to 550.84 s, lowers mean residual pressure drop to 280.69 Pa and raises quality factor to 0.044 /Pa under simulated humid dynamic filtration. The fluorine-free part is the policy-relevant part - PTFE-based antifouling treatments are the incumbent and are under PFAS restriction, so a non-fluorinated substitute that survives wet dust addresses a regulatory constraint on stationary-source PM control rather than just an engineering one. Bench-scale, simulated flue gas, no capture-efficiency-versus-size curve reported and no coating-lifetime data beyond the test campaign, so durability - the property the whole design targets - is demonstrated over hours rather than the months a baghouse runs. Tier C and upstream of this watch's usual scope; carried under mitigation because primary PM emission control determines the ambient field every downstream sensor measures.}
{10.1021/acs.est.6c08122}
% ==================================================================== TABLE
\section*{\sffamily\bfseries\color{Deep}Corpus register}
\vspace{-1.5mm}{\color{Amber}\rule{\linewidth}{1.1pt}}\vspace{2.5mm}

{\fontsize{6.7}{8.4}\selectfont
\setlength{\tabcolsep}{2.4pt}
\begin{longtable}{@{}>{\RaggedRight}p{25mm}>{\RaggedRight}p{41mm}>{\RaggedRight}p{27mm}%
>{\RaggedRight}p{23mm}>{\RaggedRight}p{18mm}>{\RaggedRight}p{15mm}@{}}
\toprule
\textsf{\textbf{First author}} & \textsf{\textbf{Focus}} & \textsf{\textbf{Journal}} &
\textsf{\textbf{Design}} & \textsf{\textbf{Metric}} & \textsf{\textbf{Region}}\\
\midrule
\endfirsthead
\toprule
\textsf{\textbf{First author}} & \textsf{\textbf{Focus}} & \textsf{\textbf{Journal}} &
\textsf{\textbf{Design}} & \textsf{\textbf{Metric}} & \textsf{\textbf{Region}}\\
\midrule
\endhead
\multicolumn{6}{@{}l@{}}{\cellcolor{Mist}\textsf{\textbf{\textcolor{Deep}{Neuro / mental health}}}}\\[0.6mm]
\href{https://doi.org/10.1001/jamanetworkopen.2026.25015}{Dales} & Wildfire-sourced PM$_{2.5}$ and ED visits for migraine and other primary headache & \textit{JAMA Netw Open} & Case-crossover & PM$_{2.5}$ (wildfire) & Canada\\
\href{https://doi.org/10.1093/ageing/afag215}{Glover} & Umbrella meta-review of urban environmental exposures and cognitive health & \textit{Age Ageing} & Umbrella review / meta-analysis & PM$_{2.5}$ & Global\\
\href{https://doi.org/10.1016/j.nsa.2026.107019}{Goossens} & In utero air pollution exposure and pre-/postnatal brain development & \textit{Neurosci Appl} & Narrative review & PM$_{2.5}$ / PM$_{10}$ / PM$_{\mathrm{coarse}}$ & Global\\
\href{https://doi.org/10.1016/j.isci.2026.116617}{Ma} & PM$_{2.5}$-induced neurotoxicity and depression: exposure risks and mechanisms & \textit{iScience} & Narrative review & PM$_{2.5}$ + components & Global\\
\href{https://doi.org/10.3760/cma.j.cn112150-20250804-00752}{Pan} & Maternal PM exposure and offspring neurodevelopment: effects and mechanisms & \textit{Zhonghua Yu Fang Yi Xue Za Zhi} & Narrative review & PM (bulk) & China\\
\addlinespace[0.9mm]
\multicolumn{6}{@{}l@{}}{\cellcolor{Mist}\textsf{\textbf{\textcolor{Teal}{Cardiovascular \& metabolic}}}}\\[0.6mm]
\href{https://doi.org/10.3760/cma.j.cn112338-20251120-00843}{Huang} & PM$_{2.5}$ and dyslipidaemia in rural Chinese children; dietary effect modification & \textit{Zhonghua Liu Xing Bing Xue Za Zhi} & Cross-sectional (multistage) & PM$_{2.5}$ & China\\
\href{https://doi.org/10.32394/pe/218966}{Kazmierska} & Air pollution and cardiovascular disease: review of current data & \textit{Przegl Epidemiol} & Narrative review & PM$_{2.5}$ / PM$_{10}$ & Global\\
\href{https://doi.org/10.1007/s00125-026-06814-2}{Xi} & Air pollution, built environment and temperature vs. prediabetes phenotypes & \textit{Diabetologia} & Prospective cohort & PM$_{2.5}$ & Germany\\
\href{https://doi.org/10.1161/JAHA.125.046917}{Yang} & Short-term PM$_{2.5}$ and hospitalisation/ED use in atrial fibrillation & \textit{J Am Heart Assoc} & Case-crossover & PM$_{2.5}$ (wildfire-influenced) & USA\\
\addlinespace[0.9mm]
\multicolumn{6}{@{}l@{}}{\cellcolor{Mist}\textsf{\textbf{\textcolor{Sky}{Reproductive \& developmental}}}}\\[0.6mm]
\href{https://doi.org/10.1016/j.jes.2025.10.043}{Chen} & Prenatal air pollution, fetal growth trajectory phenotypes and child neurodevelopment & \textit{J Environ Sci (China)} & Prospective cohort & PM$_{2.5}$ / PM$_{10}$ & China\\
\href{https://doi.org/10.1016/j.ecoenv.2026.120566}{Huang} & Joint prenatal bisphenol + air-pollutant mixture exposure and preschool IQ & \textit{Ecotoxicol Environ Saf} & Prospective cohort (mixtures) & PM$_{2.5}$ / PM$_{10}$ & China\\
\href{https://doi.org/10.1007/s40572-026-00552-8}{Peterson} & Ambient air pollution, PAHs and pubertal development: critical review & \textit{Curr Environ Health Rep} & Critical review & PM$_{2.5}$ components & Global\\
\href{https://doi.org/10.1016/j.envres.2026.125313}{Xie} & Long-term PM$_{2.5}$ and timing of pubertal development in boys and girls & \textit{Environ Res} & Prospective cohort & PM$_{2.5}$ & China\\
\addlinespace[0.9mm]
\multicolumn{6}{@{}l@{}}{\cellcolor{Mist}\textsf{\textbf{\textcolor{Sage}{Respiratory \& allergic}}}}\\[0.6mm]
\href{https://doi.org/10.1007/s00484-026-03289-z}{Bayar Muluk} & Climate change and otolaryngology: airway disease impacts and emission reduction & \textit{Int J Biometeorol} & Narrative review & PM$_{2.5}$ & Global\\
\href{https://doi.org/10.4103/mgr.MEDGASRES-D-25-00164}{Liao} & Short-term air pollution, respiratory symptoms and lung function in asthma & \textit{Med Gas Res} & Panel study & PM$_{2.5}$ / PM$_{10}$ & China\\
\href{https://doi.org/10.3390/epidemiologia7040097}{Sanchez de Prada} & Ambient PM and asthma-coded admissions in children 0-3 y, Spain & \textit{Epidemiologia} & Ecological panel & PM$_{2.5}$ / PM$_{10}$ & Spain\\
\href{https://doi.org/10.1016/j.jhazmat.2026.143076}{Yang} & Allergenic vs. toxic risk of airborne seed hairs compared with ambient PM$_{2.5}$ & \textit{J Hazard Mater} & Field measurement & PM$_{2.5}$ + bioaerosol & China\\
\addlinespace[0.9mm]
\multicolumn{6}{@{}l@{}}{\cellcolor{Mist}\textsf{\textbf{\textcolor{Violet}{Mechanistic toxicology}}}}\\[0.6mm]
\href{https://doi.org/10.1016/j.jes.2025.12.074}{Ding} & PM$_{2.5}$ promotes vascular calcification via miR-27b-5p/HES1 signalling & \textit{J Environ Sci (China)} & In vitro + organoid + murine & PM$_{2.5}$ & China\\
\href{https://doi.org/10.1007/s00360-026-01698-5}{Maraschi} & Osmoregulatory disruption and metal bioaccumulation from settleable atmospheric PM & \textit{J Comp Physiol B} & Controlled exposure (ecotox) & Settleable APM & Brazil\\
\href{https://doi.org/10.1016/j.jes.2025.11.018}{Wu} & Macrophage NLRP3 inflammasome in real-ambient PM-induced hepatic glucose dysregulation & \textit{J Environ Sci (China)} & Real-ambient murine + co-culture & Real-ambient PM & China\\
\addlinespace[0.9mm]
\multicolumn{6}{@{}l@{}}{\cellcolor{Mist}\textsf{\textbf{\textcolor{Clay}{Occupational \& indoor}}}}\\[0.6mm]
\href{https://doi.org/10.3389/fmedt.2026.1796828}{Chen} & Nested suction instrument to cut intraoperative surgical-smoke PM$_{2.5}$ & \textit{Front Med Technol} & Bench / device experiment & PM$_{2.5}$ (surgical smoke) & China\\
\href{https://doi.org/10.1093/annweh/wxag055}{Cole} & Determinants of respirable crystalline silica exposure in Australian tunnelling & \textit{Ann Work Expo Health} & Observational exposure study & Respirable dust / RCS & Australia\\
\href{https://doi.org/10.1016/j.jhazmat.2026.143081}{Ding} & Stage-resolved hair Zn-isotope dynamics as a firefighting combustion-exposure tracer & \textit{J Hazard Mater} & Repeated-measures biomarker & Combustion-derived PM & China\\
\href{https://doi.org/10.3390/jox16040119}{Ramsden} & Xenobiotic hazards in aircraft cabin air & \textit{J Xenobiot} & Review + risk economics & Ultrafine particles & Global\\
\addlinespace[0.9mm]
\multicolumn{6}{@{}l@{}}{\cellcolor{Mist}\textsf{\textbf{\textcolor{Amber}{Sensing, forecasting \& instrumentation}}}}\\[0.6mm]
\href{https://doi.org/10.1016/j.jenvman.2026.130515}{Shi} & Adaptive spatiotemporal graph transformer for multi-site multi-step PM$_{2.5}$ forecasting & \textit{J Environ Manage} & Deep-learning model & PM$_{2.5}$ & China / India\\
\href{https://doi.org/10.3390/bios16070380}{Tsuchiya} & Seasonal and indoor-outdoor characteristics of fluorescent aerosol particles in offices & \textit{Biosensors} & Real-time sensor campaign & Bioaerosol / coarse & Japan\\
\addlinespace[0.9mm]
\multicolumn{6}{@{}l@{}}{\cellcolor{Mist}\textsf{\textbf{\textcolor{Amber}{Exposure assessment \& modelling}}}}\\[0.6mm]
\href{https://doi.org/10.1029/2025GH001621}{Solorzano-Araque} & Multiscale exposure modelling: LOTOS-EUROS + mobility surveys + particle toxicity & \textit{GeoHealth} & CTM + mobility integration & PM$_{2.5}$ / PM$_{10}$ & Colombia\\
\addlinespace[0.9mm]
\multicolumn{6}{@{}l@{}}{\cellcolor{Mist}\textsf{\textbf{\textcolor{Coral}{Burden, policy \& mitigation}}}}\\[0.6mm]
\href{https://doi.org/10.1002/hsr2.72886}{Karamzad} & Ischaemic heart disease burden in MENA 1990-2021; PM attributable fraction 28.2\% & \textit{Health Sci Rep} & GBD secondary analysis & Ambient PM & MENA\\
\href{https://doi.org/10.1371/journal.pone.0347989}{Rees} & HEPA filtration to reduce respiratory infection in care homes & \textit{PLoS One} & RCT process evaluation & Filtered PM & UK\\
\href{https://doi.org/10.1016/j.jes.2026.01.009}{Shan} & Policy-driven emission reductions and post-pandemic rebound, high-resolution inventories & \textit{J Environ Sci (China)} & Emission inventory & PM emissions & China\\
\addlinespace[0.9mm]
\multicolumn{6}{@{}l@{}}{\cellcolor{Mist}\textsf{\textbf{\textcolor{Slate}{Other clinical endpoints}}}}\\[0.6mm]
\href{https://doi.org/10.3390/dj14070440}{Lupita} & PM$_{2.5}$, preventive-care capacity and caries/edentulism burden across the EU & \textit{Dent J} & Ecological panel & PM$_{2.5}$ & EU-27\\
\href{https://doi.org/10.1016/j.jdin.2026.06.008}{Park} & Long-term PM exposure and risk of androgenetic alopecia & \textit{JAAD Int} & Nationwide cohort & PM$_{2.5}$ / PM$_{10}$ & South Korea\\
\href{https://doi.org/10.3171/2026.3.PEDS2615}{Sun} & Environmental exposure and area deprivation vs. paediatric hydrocephalus outcomes & \textit{J Neurosurg Pediatr} & Registry cohort + geospatial & PM$_{2.5}$ (satellite-derived) & USA\\
\addlinespace[0.9mm]
\bottomrule
\end{longtable}}


\vspace{4mm}
\begin{keybox}
{\sffamily\bfseries\footnotesize\color{Deep}Provenance and method}\\[1.2mm]
{\fontsize{7.2}{8.8}\selectfont
Entry window \textbf{21 September 2026}, single day, continuous with the previous issue and
leaving no gap. Legs run: \textbf{PubMed} E-utilities on the health and instrumentation
axes as separate queries (\textbf{6} health, \textbf{0} sensing --- the sensing leg returned
nothing for the \textbf{second consecutive issue}); \textbf{Europe PMC}
\texttt{CREATION\_DATE} (\textbf{5}); \textbf{Crossref by ISSN} (\textbf{36}).
\textbf{OpenAlex} returned HTTP 429 and \textbf{arXiv} HTTP 406, both known-broken legs
carried in the source-status table. No Europe PMC records were dropped as retro-index
artefacts this window. The \textbf{47} raw records deduplicated to \textbf{46}, of which
\textbf{1} was already carried in a previous issue, leaving \textbf{45} fresh candidates of
which \textbf{18 were admitted} and 27 rejected, every one logged with a reason. The
rejected set is dominated by two journals whose scope overlaps this watch only partially:
seven \textit{Environmental Science \& Technology} records on water and sediment chemistry,
and four \textit{Measurement Science and Technology} records on interferometry, acoustics,
EEG separation and fringe projection --- instrument papers that share the shelf but not the
measurand. One record was rejected for \textbf{carrying no DOI and no PMID} (a murine
exposome cardiovascular study), the sole DOI-gate casualty this window. Following the
20 September finding that a publisher had registered a DOI within a day of a
DOI-gate rejection, the recent DOI-less rejects were re-checked before screening; none had
acquired one.
\textbf{Five records are carried on metadata alone} and assert no finding. Recovery was
attempted for all five through PubMed \texttt{efetch}, Crossref and Europe PMC, and for two
of them additionally through the publisher landing page and a Consensus sweep. The tier-A
record of this issue was itself recovered that way --- it deposited with \textbf{no abstract
and no PMID}, and its effect sizes come from the publisher landing page, which is the only
reason this issue has a forest plot at all. Semantic Scholar was not retried, having been
established as dead for same-day deposits on 19 September. All \textbf{18} DOIs resolved
before build (\textbf{0 fail, 0 warn}). One record
(\texttt{10.1021/acs.est.6c08122}) tripped a \textsc{not-in-crossref} warning on the first
build pass and cleared on the second within the same run --- an ACS deposit indexing while
the issue was being assembled, which is the same one-day registration lag that recovered a
DOI-gate reject on 20 September.
Subtopic, design, particle-metric and geography labels are assigned from title, abstract,
keywords and MeSH; assignment is single-label by dominant contribution and is a judgement,
not a controlled vocabulary. \textbf{Thirteen new design labels} were registered in
\texttt{plots.py} in the same edit that wrote them onto records --- the largest single-issue
batch since 23 August, and none of them a new \textit{kind} of study. The five abstract-less
records were deliberately \textbf{not} given bespoke labels; they use the canonical
\textit{Metadata only (no abstract)} string so they pool into one honest slice rather than
fragmenting the architecture figure into five singletons. Two further defects were fixed in
this pass: the \textbf{forest plot's latent ratio-scale assumption} (see \textit{What changed
today}) and \textbf{two bench geography labels} that would otherwise have inflated the
\textit{Global / multi-region} bar. Effect estimates are transcribed verbatim and use
study-specific exposure contrasts --- see the forest-plot caption. Abstracts remain the
copyright of their respective publishers.\par}
\end{keybox}

\end{document}
